%% Lecture 20 -- Waves, Lorenz gauge, and retarded Green functions
\section{Lecture 20 -- Waves and Retarded Green Functions}\label{sec:lec20}

\begin{center}
\begin{tabular}{ll}
  \textbf{Date:}\quad 14/06/2026 & \\
\end{tabular}
\end{center}
\hrule\vspace{1.5em}

\subsection*{Overview}

We now leave magnetostatics and return to the full time-dependent Maxwell equations. The
central object is still the four-potential $A^\mu=(\phi,\Avec)$ and the field tensor
$F^{\mu\nu}=\partial^\mu A^\nu-\partial^\nu A^\mu$, but the source $J^\mu$ may now depend
on time. This changes the mathematical problem from Poisson's equation to a wave equation:
information no longer fills all space instantaneously, but propagates on the light cone.

The lecture has three connected goals. First, we use gauge freedom to impose the Lorenz
condition $\partial_\mu A^\mu=0$ and reduce the inhomogeneous Maxwell equations to four
identical sourced wave equations. Second, we formulate the Green-function solution of this
wave equation and check that current conservation automatically keeps the Lorenz gauge
consistent. Third, we solve the scalar Green-function equation by Fourier transforming only
in time. This converts the wave equation into the Helmholtz equation, whose spherically
symmetric Green functions are outgoing and incoming spherical waves. Choosing the outgoing
one is the physical 
\emph{retarded} boundary condition.

The conceptual point is causality. The differential equation by itself admits both
retarded and advanced Green functions: a pulse can mathematically propagate toward the
future or toward the past. Physics enters through the boundary condition that sources
produce fields after they act, not before.

%------------------------------------------------------------
\subsection{Maxwell's Equations as a Wave Equation for $A^\mu$}\label{subsec:lec20-maxwell-wave}

\subsubsection*{Conceptual Background}
In electrostatics and magnetostatics we solved elliptic equations such as
$\nabla^2\phi=-4\pi\rho$ or $-\nabla^2\Avec=4\pi\Jvec/c$. Those equations describe
instantaneous constraints on a single time slice. Once $\rho$ and $\Jvec$ depend on time,
the correct relativistic object is the four-current $J^\mu=(c\rho,\Jvec)$, and charge
conservation becomes the covariant continuity equation
\begin{equation}\label{eq:lec20-continuity}
  \boxed{\partial_\mu J^\mu=0.}
\end{equation}
\textit{Significance:} The source of Maxwell's equations is physically admissible only if
charge is neither created nor destroyed.

The inhomogeneous Maxwell equations are
\begin{equation}\label{eq:lec20-maxwell-covariant}
  \boxed{\partial_\mu F^{\mu\nu}=\frac{4\pi}{c}J^\nu.}
\end{equation}
\textit{Significance:} This single tensor equation contains Gauss's law and the
Ampere--Maxwell law, while the homogeneous equations are already built into
$F_{\mu\nu}=\partial_\mu A_\nu-\partial_\nu A_\mu$.

Insert
\begin{equation}\label{eq:lec20-F-from-A}
  F^{\mu\nu}=\partial^\mu A^\nu-\partial^\nu A^\mu
\end{equation}
into \eqref{eq:lec20-maxwell-covariant}. Since partial derivatives commute,
\begin{align*}
  \partial_\mu F^{\mu\nu}
  &=\partial_\mu\qty(\partial^\mu A^\nu-\partial^\nu A^\mu) \\
  &=\partial_\mu\partial^\mu A^\nu-\partial_\mu\partial^\nu A^\mu \\
  &=\partial_\mu\partial^\mu A^\nu-\partial^\nu\qty(\partial_\mu A^\mu).
\end{align*}
Therefore the exact equation for the potential is
\begin{equation}\label{eq:lec20-A-before-gauge}
  \boxed{\partial_\mu\partial^\mu A^\nu
  -\partial^\nu\qty(\partial_\mu A^\mu)=\frac{4\pi}{c}J^\nu.}
\end{equation}
\textit{Significance:} The four components of $A^\mu$ are coupled only through the gauge-dependent scalar $\partial_\mu A^\mu$.

\subsubsection*{The d'Alembertian}
With the metric $\eta_{\mu\nu}=\mathrm{diag}(+1,-1,-1,-1)$ and $x^0=ct$,
\begin{align*}
  \partial_\mu\partial^\mu
  &=\eta^{\mu\nu}\partial_\mu\partial_\nu
    =\eta^{00}\partial_0^2+\sum_{i=1}^{3}\eta^{ii}\partial_i^2 \\
  &=\frac{1}{c^2}\frac{\partial^2}{\partial t^2}-\nabla^2.
\end{align*}
We denote this operator by
\begin{equation}\label{eq:lec20-dalembertian}
  \boxed{\partial_\mu\partial^\mu\equiv \Box
  =\frac{1}{c^2}\frac{\partial^2}{\partial t^2}-\nabla^2.}
\end{equation}
\textit{Significance:} $\Box$ is the Lorentz-invariant wave operator; replacing $\nabla^2$ by $-\Box$ is the mathematical transition from statics to radiation.

%------------------------------------------------------------
\subsection{Gauge Fixing: The Lorenz Gauge}\label{subsec:lec20-lorenz-gauge}

\subsubsection*{Physical Motivation}
The potential is not unique: the transformation
\begin{equation}\label{eq:lec20-gauge-transformation}
  \boxed{A^\mu\longrightarrow \widetilde A^\mu=A^\mu+\partial^\mu f}
\end{equation}
\textit{Significance:} Gauge freedom changes the potential representation but leaves the measurable fields $\bE$ and $\bB$ unchanged.

We use this freedom to choose a representative satisfying the Lorenz gauge
\begin{equation}\label{eq:lec20-lorenz-condition}
  \boxed{\partial_\mu A^\mu=0.}
\end{equation}
\textit{Significance:} The Lorenz gauge is the covariant analogue of the Coulomb gauge and turns Maxwell's equations into decoupled wave equations.

\subsubsection*{Existence of the Lorenz Gauge}
Suppose some potential $A^\mu$ does not obey \eqref{eq:lec20-lorenz-condition}. Define the
new potential by \eqref{eq:lec20-gauge-transformation}. Its four-divergence is
\begin{align*}
  \partial_\mu\widetilde A^\mu
  &=\partial_\mu A^\mu+\partial_\mu\partial^\mu f \\
  &=\partial_\mu A^\mu+\Box f.
\end{align*}
Demanding $\partial_\mu\widetilde A^\mu=0$ gives the sourced wave equation
\begin{equation}\label{eq:lec20-gauge-f-equation}
  \boxed{\Box f=-\partial_\mu A^\mu.}
\end{equation}
\textit{Significance:} Imposing Lorenz gauge is legitimate because it reduces to solving the same kind of wave equation that we are about to solve for the fields.

In the Lorenz gauge the second term of \eqref{eq:lec20-A-before-gauge} vanishes, and
Maxwell's equations become
\begin{equation}\label{eq:lec20-wave-A}
  \boxed{\Box A^\nu=\frac{4\pi}{c}J^\nu.}
\end{equation}
\textit{Significance:} Each component of the four-potential propagates as a wave sourced by the corresponding component of the four-current.

\nt{Pitfall: the name is \emph{Lorenz} gauge, after Ludvig Lorenz, not Lorentz gauge after Hendrik Lorentz. The condition is Lorentz covariant, but the historical name has no ``t''.}

\subsubsection*{Method: fixing Lorenz gauge in practice}
\begin{enumerate}
  \item Compute the scalar $S\equiv\partial_\mu A^\mu$ for the potential you are given.
  \item Solve the wave equation $\Box f=-S$ with the boundary condition appropriate to the problem.
  \item Replace $A^\mu$ by $\widetilde A^\mu=A^\mu+\partial^\mu f$.
  \item Check explicitly that $\partial_\mu\widetilde A^\mu=S+\Box f=0$.
\end{enumerate}

\ex[ex:lec20-gauge-fix]{Worked Example: Building a Lorenz-Gauge Representative}{
Suppose a potential has divergence
\begin{equation*}
  \partial_\mu A^\mu=S(x),
\end{equation*}
where $S(x)$ is known. Let $G(x-y)$ be any Green function satisfying
$\Box_xG(x-y)=\delta^{(4)}(x-y)$. Choose
\begin{equation*}
  f(x)=-\int d^4y\,G(x-y)S(y).
\end{equation*}
Then
\begin{align*}
  \Box_x f(x)
  &=-\int d^4y\,\Box_xG(x-y)S(y) \\
  &=-\int d^4y\,\delta^{(4)}(x-y)S(y) \\
  &=-S(x).
\end{align*}
Therefore
\begin{equation*}
  \partial_\mu\widetilde A^\mu=\partial_\mu A^\mu+\Box f=S-S=0.
\end{equation*}
This explicitly constructs a Lorenz-gauge potential from the original one.
}

%------------------------------------------------------------
\subsection{Green-Function Solution of the Sourced Wave Equation}\label{subsec:lec20-green-solution}

\subsubsection*{Mathematical Setup}
Equation \eqref{eq:lec20-wave-A} is linear and inhomogeneous. Therefore its general
solution has the form
\begin{equation}\label{eq:lec20-general-solution}
  \boxed{\begin{aligned}
    A^\nu&=A^\nu_{\rm p}+A^\nu_{\rm h},\\
    \Box A^\nu_{\rm p}&=\frac{4\pi}{c}J^\nu,\qquad
    \Box A^\nu_{\rm h}=0.
  \end{aligned}}
\end{equation}
\textit{Significance:} The particular solution is the field generated by the sources, while the homogeneous solution is incoming free radiation supplied from outside the source region.

A scalar Green function for the wave operator is defined by
\begin{equation}\label{eq:lec20-green-definition}
  \boxed{\Box_xG(x-y)=\delta^{(4)}(x-y),}
\end{equation}
where $x^\mu=(ct,\br)$ and $y^\mu=(ct',\br')$ are spacetime points.
\textit{Significance:} $G(x-y)$ is the response to a unit spacetime impulse located at $y$.

Using \eqref{eq:lec20-green-definition}, a particular solution is
\begin{equation}\label{eq:lec20-A-particular}
  \boxed{A^\nu_{\rm p}(x)=\frac{4\pi}{c}\int d^4y\,G(x-y)J^\nu(y).}
\end{equation}
\textit{Significance:} Once the response to a point impulse is known, the response to an arbitrary source is obtained by superposition over spacetime.

To verify \eqref{eq:lec20-A-particular}, apply $\Box_x$:
\begin{align*}
  \Box_x A^\nu_{\rm p}(x)
  &=\frac{4\pi}{c}\int d^4y\,\Box_xG(x-y)J^\nu(y) \\
  &=\frac{4\pi}{c}\int d^4y\,\delta^{(4)}(x-y)J^\nu(y) \\
  &=\frac{4\pi}{c}J^\nu(x),
\end{align*}
which is exactly \eqref{eq:lec20-wave-A}.

\subsubsection*{Consistency with the Lorenz Gauge}
We must also check that the Green-function solution obeys $\partial_\nu A^\nu=0$ if the
source is conserved. Differentiate \eqref{eq:lec20-A-particular}:
\begin{align*}
  \partial^x_\nu A^\nu_{\rm p}(x)
  &=\frac{4\pi}{c}\int d^4y\,\partial^x_\nu G(x-y)J^\nu(y).
\end{align*}
Because $G$ depends only on $x-y$,
\begin{equation*}
  \partial^x_\nu G(x-y)=-\partial^y_\nu G(x-y).
\end{equation*}
Therefore
\begin{align*}
  \partial^x_\nu A^\nu_{\rm p}(x)
  &=-\frac{4\pi}{c}\int d^4y\,\partial^y_\nu G(x-y)J^\nu(y) \\
  &=-\frac{4\pi}{c}\int d^4y\,\partial^y_\nu\qty(GJ^\nu)
    +\frac{4\pi}{c}\int d^4y\,G(x-y)\partial^y_\nu J^\nu(y).
\end{align*}
For localized sources and Green functions obeying physical boundary conditions, the total
four-divergence gives a boundary term at spatial and temporal infinity and vanishes. The
second term vanishes by \eqref{eq:lec20-continuity}. Thus
\begin{equation}\label{eq:lec20-lorenz-preserved}
  \boxed{\partial_\nu A^\nu_{\rm p}=0\qquad(\partial_\nu J^\nu=0).}
\end{equation}
\textit{Significance:} Charge conservation is precisely what makes the Lorenz-gauge wave equation self-consistent.

\nt{Pitfall: the Green function itself is only an auxiliary object and its source
$\delta^{(4)}(x-y)$ is not a conserved electromagnetic current. Conservation is required
for the physical $J^\mu$ inserted in \eqref{eq:lec20-A-particular}, not for the mathematical
impulse defining $G$.}

\subsubsection*{Method: solving a sourced wave equation with a Green function}
\begin{enumerate}
  \item Choose boundary conditions first: retarded, advanced, incoming radiation, outgoing radiation, or a finite-domain condition.
  \item Solve $\Box G=\delta^{(4)}$ with those boundary conditions.
  \item Convolve: $A^\nu_{\rm p}=(4\pi/c)G*J^\nu$.
  \item Add any desired homogeneous solution $A^\nu_{\rm h}$ if the physical problem includes externally incident radiation.
  \item Verify the Lorenz condition using $\partial_\mu J^\mu=0$ and the boundary behavior.
\end{enumerate}

%------------------------------------------------------------
\subsection{Solving the Wave Green Function}\label{subsec:lec20-solving-green}

\subsubsection*{Mathematical Setup}
Translation invariance lets us place the impulse at the origin, $y=0$, and write
$G(x-y)=G(x)$. The equation becomes
\begin{equation}\label{eq:lec20-green-origin}
  \qty(\frac{\partial^2}{\partial (x^0)^2}-\nabla^2)G(x^0,\br)=\delta(x^0)\delta^{(3)}(\br),
\end{equation}
where $x^0=ct$. We Fourier transform only in time:
\begin{equation}\label{eq:lec20-fourier-def}
  \widetilde G(k_0,\br)=\int_{-\infty}^{\infty}dx^0\,e^{-ik_0x^0}G(x^0,\br),
\end{equation}
with inverse
\begin{equation}\label{eq:lec20-fourier-inverse}
  G(x^0,\br)=\int_{-\infty}^{\infty}\frac{dk_0}{2\pi}\,e^{ik_0x^0}\widetilde G(k_0,\br).
\end{equation}
The time derivative term is integrated by parts twice. Assuming the boundary terms vanish,
\begin{align*}
  \int dx^0\,e^{-ik_0x^0}\frac{\partial^2G}{\partial (x^0)^2}
  &=\int dx^0\,\frac{\partial^2}{\partial (x^0)^2}\qty(e^{-ik_0x^0})G \\
  &=(-ik_0)^2\widetilde G=-k_0^2\widetilde G.
\end{align*}
The Fourier transform of $\delta(x^0)$ is $1$. Hence \eqref{eq:lec20-green-origin}
becomes the Helmholtz Green equation
\begin{equation}\label{eq:lec20-helmholtz}
  \boxed{\qty(\nabla^2+k_0^2)\widetilde G(k_0,\br)=-\delta^{(3)}(\br).}
\end{equation}
\textit{Significance:} At fixed frequency $k_0$, the wave Green function is a spherical Helmholtz wave sourced at the origin.

\subsubsection*{Rotational Symmetry and the Radial Equation}
The source $\delta^{(3)}(\br)$ is rotationally invariant, so we seek a rotationally
invariant solution $\widetilde G(k_0,\br)=\widetilde G(k_0,r)$ with $r=|\br|$. For a radial
function,
\begin{equation}\label{eq:lec20-radial-laplacian}
  \nabla^2\widetilde G=\frac{1}{r}\frac{d^2}{dr^2}\qty(r\widetilde G).
\end{equation}
For $r\neq0$ the delta function vanishes, and \eqref{eq:lec20-helmholtz} gives
\begin{equation}\label{eq:lec20-radial-homogeneous}
  \frac{d^2}{dr^2}\qty(r\widetilde G)+k_0^2\qty(r\widetilde G)=0.
\end{equation}
Thus
\begin{equation}\label{eq:lec20-Gtilde-general}
  \boxed{\widetilde G(k_0,r)=A\frac{e^{-ik_0r}}{r}+B\frac{e^{ik_0r}}{r},\qquad r\neq0.}
\end{equation}
\textit{Significance:} The frequency-domain point response is a superposition of two spherical waves, one associated with propagation toward the future and one toward the past after transforming back to time.

The singularity at $r=0$ fixes only $A+B$. Use the distributional identity
\begin{equation}\label{eq:lec20-laplacian-one-over-r}
  \nabla^2\frac{1}{r}=-4\pi\delta^{(3)}(\br).
\end{equation}
Near $r=0$, $e^{\pm ik_0r}=1+O(r)$, so the singular part of \eqref{eq:lec20-Gtilde-general} is
$(A+B)/r$. Therefore
\begin{equation}\label{eq:lec20-AplusB}
  \boxed{A+B=\frac{1}{4\pi}.}
\end{equation}
\textit{Significance:} The local Coulomb-like singularity fixes the normalization, while the split between $A$ and $B$ remains a boundary-condition choice.

\nt{Pitfall: do not multiply radial delta functions by powers of $r$ naively. The identity
$f(\br)\delta^{(3)}(\br)=f(0)\delta^{(3)}(\br)$ requires $f$ to be regular at the origin;
objects such as $r\delta^{(3)}(\br)$ inside spherical-coordinate measures must be handled as distributions, not by informal substitution.}

\subsubsection*{Fourier Transform Back to Time}
Using \eqref{eq:lec20-fourier-inverse},
\begin{align*}
  G(x^0,r)
  &=\int\frac{dk_0}{2\pi}e^{ik_0x^0}
    \qty(A\frac{e^{-ik_0r}}{r}+B\frac{e^{ik_0r}}{r}) \\
  &=\frac{A}{r}\int\frac{dk_0}{2\pi}e^{ik_0(x^0-r)}
    +\frac{B}{r}\int\frac{dk_0}{2\pi}e^{ik_0(x^0+r)} \\
  &=\frac{A}{r}\delta(x^0-r)+\frac{B}{r}\delta(x^0+r).
\end{align*}
With \eqref{eq:lec20-AplusB}, the two elementary causal choices are
\begin{equation}\label{eq:lec20-ret-adv-green}
  \boxed{G_{\rm ret}(x)=\frac{\delta(x^0-r)}{4\pi r},\qquad
  G_{\rm adv}(x)=\frac{\delta(x^0+r)}{4\pi r}.}
\end{equation}
\textit{Significance:} $G_{\rm ret}$ is nonzero only on the future light cone of the source event, while $G_{\rm adv}$ is nonzero only on its past light cone.

Restoring the source point $y=(ct',\br')$ and writing $R=|\br-\br'|$ gives
\begin{equation}\label{eq:lec20-ret-green-restored}
  \boxed{G_{\rm ret}(x-y)=\frac{\delta\qty(c(t-t')-R)}{4\pi R},
  \qquad R=|\br-\br'|.}
\end{equation}
\textit{Significance:} A source element at $\br'$ affects the observation point $\br$ only after the light-travel time $R/c$.

\begin{figure}[h]
\centering
\begin{tikzpicture}[>=Stealth, scale=1.0]
  \draw[->] (-0.2,0) -- (6.1,0) node[right] {$x^1$};
  \draw[->] (0,-0.3) -- (0,4.2) node[above] {$x^0=ct$};
  \coordinate (O) at (0,0);
  \draw[fill=black] (O) circle (1.5pt) node[below left] {source event $y$};
  \draw[blue!70!black, thick] (O) -- (3.4,3.4) node[above right] {future light cone};
  \draw[blue!70!black, thick] (O) -- (-0.9,0.9);
  \draw[red!70!black, dashed, thick] (O) -- (3.0,-3.0);
  \draw[red!70!black, dashed, thick] (O) -- (-0.9,-0.9) node[left] {past cone};
  \draw[->, thick, orange!80!black] (0,0) -- (2.2,2.2);
  \node[orange!80!black, right] at (1.35,1.65) {$R=c(t-t')$};
  \draw[dotted] (2.2,0) -- (2.2,2.2);
  \node[below] at (2.2,0) {$R$};
  \node[blue!70!black] at (4.5,2.8) {$G_{\rm ret}\neq0$};
\end{tikzpicture}
\caption{Retarded support in a spacetime diagram. The impulse at $y$ produces a response only on its future light cone, $c(t-t')=|\br-\br'|$. The advanced Green function would instead live on the past cone.}
\label{fig:lec20-light-cone}
\end{figure}

\subsubsection*{Method: deriving the retarded Green function}
\begin{enumerate}
  \item Use translation invariance to set the impulse at the origin.
  \item Fourier transform in $x^0$ to convert $\Box G=\delta^{(4)}$ into the Helmholtz equation.
  \item Use rotational symmetry to reduce the Helmholtz equation to an ordinary radial equation.
  \item Solve for $r\neq0$: $\widetilde G=(Ae^{-ik_0r}+Be^{ik_0r})/r$.
  \item Fix $A+B=1/(4\pi)$ from the $1/r$ singularity.
  \item Choose the retarded boundary condition $B=0$ so the response vanishes before the source acts.
\end{enumerate}

\ex[ex:lec20-static-limit]{Worked Example: Retarded Potential of a Static Point Charge}{
Take a charge $q$ at rest at the origin. Its four-current is
\begin{equation*}
  J^0(y)=c q\delta^{(3)}(\br'),\qquad \Jvec(y)=\bvec{0}.
\end{equation*}
Using \eqref{eq:lec20-A-particular} with \eqref{eq:lec20-ret-green-restored},
\begin{align*}
  \phi(t,\br)=A^0(t,\br)
  &=\frac{4\pi}{c}\int cdt'\dd[3]{r}'\,
    \frac{\delta\qty(c(t-t')-|\br-\br'|)}{4\pi|\br-\br'|}
    c q\delta^{(3)}(\br') \\
  &=4\pi c q\int dt'\,
    \frac{\delta\qty(c(t-t')-r)}{4\pi r} \\
  &=\frac{c q}{r}\int dt'\,\frac{1}{c}\delta\qty(t'-t+\frac{r}{c})
    =\frac{q}{r}.
\end{align*}
The spatial potential is
\begin{equation*}
  \Avec(t,\br)=\frac{4\pi}{c}\int d^4y\,G_{\rm ret}(x-y)\Jvec(y)=\bvec{0}.
\end{equation*}
Hence the retarded formalism reproduces the Coulomb result for a time-independent source:
\begin{equation*}
  \phi=\frac{q}{r},\qquad \Avec=\bvec{0},\qquad \bE=\frac{q\br}{r^3},\qquad \bB=\bvec{0}.
\end{equation*}
}

%------------------------------------------------------------
\subsection{Retarded Potentials}\label{subsec:lec20-retarded-potentials}

\subsubsection*{Physical Motivation}
The retarded Green function is the physically standard choice because it implements the
boundary condition ``no field before the source acts.'' Insert
\eqref{eq:lec20-ret-green-restored} into \eqref{eq:lec20-A-particular}. Since
$d^4y=dy^0d^3r'=c\dd{t}'\dd[3]{r}'$,
\begin{align*}
  A^\nu(t,\br)
  &=\frac{4\pi}{c}\int cdt'\dd[3]{r}'\,
    \frac{\delta\qty(c(t-t')-R)}{4\pi R}J^\nu(t',\br') \\
  &=4\pi\int dt'\dd[3]{r}'\,
    \frac{\delta\qty(c(t-t')-R)}{4\pi R}J^\nu(t',\br').
\end{align*}
Use
\begin{equation*}
  \delta\qty(c(t-t')-R)=\frac{1}{c}\delta\qty(t'-t+\frac{R}{c}),
\end{equation*}
so the $t'$ integral evaluates the source at the retarded time
\begin{equation}\label{eq:lec20-retarded-time}
  \boxed{t_{\rm ret}=t-\frac{|\br-\br'|}{c}.}
\end{equation}
\textit{Significance:} The field observed at $(t,\br)$ depends on what the source at $\br'$ was doing when light left it.

Because $J^0=c\rho$ and $J^i=J_i$ as spatial components of the current vector, the standard
retarded potentials are
\begin{equation}\label{eq:lec20-retarded-potentials}
  \boxed{\phi(t,\br)=\int d^3r'\,\frac{\rho(t_{\rm ret},\br')}{|\br-\br'|},\qquad
  \Avec(t,\br)=\frac{1}{c}\int d^3r'\,\frac{\Jvec(t_{\rm ret},\br')}{|\br-\br'|}.}
\end{equation}
\textit{Significance:} Time-dependent electrodynamics is the static Coulomb/Biot--Savart superposition evaluated at retarded time, with radiation effects hidden in the time dependence of the sources.

\nt{Pitfall: $R=|\br-\br'|$ depends on the integration point $\br'$. When differentiating
retarded potentials to compute $\bE$ and $\bB$, derivatives act both on $1/R$ and on the
retarded source through $t_{\rm ret}=t-R/c$. Forgetting this is the standard way to miss
radiation terms.}

%------------------------------------------------------------
\subsection{Summary}\label{subsec:lec20-summary}

\begin{itemize}
  \item The full time-dependent Maxwell equations in Lorenz gauge reduce to
        $\Box A^\nu=4\pi J^\nu/c$, with $\Box=c^{-2}\partial_t^2-\nabla^2$.
  \item Lorenz gauge can always be reached by solving $\Box f=-\partial_\mu A^\mu$ and
        transforming $A^\mu\to A^\mu+\partial^\mu f$.
  \item A wave Green function obeys $\Box G(x-y)=\delta^{(4)}(x-y)$, giving
        $A^\nu_{\rm p}=(4\pi/c)\int d^4y\,G(x-y)J^\nu(y)$.
  \item Current conservation $\partial_\mu J^\mu=0$ guarantees that the Green-function
        solution respects the Lorenz condition, up to boundary terms at infinity.
  \item Fourier transforming in time turns the Green equation into
        $(\nabla^2+k_0^2)\widetilde G=-\delta^{(3)}(\br)$, whose radial solutions are
        $(Ae^{-ik_0r}+Be^{ik_0r})/r$ with $A+B=1/(4\pi)$.
  \item Choosing $B=0$ gives the retarded Green function
        $G_{\rm ret}=\delta(x^0-r)/(4\pi r)$, hence the retarded potentials
        $\phi=\int \rho(t-R/c,\br')/R\dd[3]{r}'$ and
        $\Avec=c^{-1}\int \Jvec(t-R/c,\br')/R\dd[3]{r}'$.
\end{itemize}

\subsection*{Further Reading}
\begin{itemize}
  \item J.\,D. Jackson, \textit{Classical Electrodynamics}, 3rd ed.: Chapter 6 (time-dependent Maxwell equations, gauges, Green functions, and retarded potentials).
  \item D.\,J. Griffiths, \textit{Introduction to Electrodynamics}, 4th ed.: Chapter 10 (electrodynamic potentials and retarded solutions).
  \item L.\,D. Landau \& E.\,M. Lifshitz, \textit{The Classical Theory of Fields}, Vol.\ 2: Chapter 8 (the electromagnetic field equations and retarded potentials).
  \item David Tong, \textit{Lectures on Electromagnetism}, Chapter 6 notes -- \url{https://www.damtp.cam.ac.uk/user/tong/em.html}
  \item MIT OpenCourseWare 8.07, \textit{Electromagnetism II} -- \url{https://ocw.mit.edu/courses/8-07-electromagnetism-ii-fall-2012/}
\end{itemize}
